\section*{Appendix D: Conventions and Notation}
\addcontentsline{toc}{section}{Appendix D: Conventions and Notation}

\section{Indices and Summation}

Throughout this monograph, lower case Latin indices $i,j,k,\dots$ range over spatial components $1,\dots,n$, while Greek indices $\mu,\nu,\rho,\sigma$ range over spacetime indices $0,\dots,n$.  The Einstein summation convention over repeated indices is always assumed unless stated otherwise.

Boldface symbols such as $\mathbf v$ denote spatial vector fields, whereas non-bold symbols $v^\mu$ denote spacetime components.

\section{Metrics and Signatures}

Spacetime metrics $g_{\mu\nu}$ have Lorentzian signature $(-,+,\ldots,+)$ unless otherwise stated.  Raising and lowering of indices is performed with $g_{\mu\nu}$ and its inverse $g^{\mu\nu}$.

The metric used in classical background comparisons may differ from the effective geometric structures arising in lamphrodynamics; the latter are derived from the triple $(\Phi,\mathbf v,S)$ rather than postulated \emph{a priori}.

\section{Differential Operators}

The gradient operator $\nabla$ denotes the Levi--Civita covariant derivative unless the context dictates a gauge or derived connection. The d'Alembertian operator is
[
\square = \nabla^\mu\nabla_\mu.
]

The divergence of a vector field is $\nabla\cdot\mathbf v$, and the Laplacian on scalar functions is $\Delta=\nabla^i\nabla_i$.  In derived settings, the operator $\Delta$ refers to the BV Laplacian unless explicitly disambiguated.

\section{Lamphrodynamic Fields}

The symbol $\Phi$ denotes the scalar potential field.  The bold symbol $\mathbf v$ denotes the spatial vector field of lamphrodynamic flow, and $S$ denotes the entropy density.  The triple $(\Phi,\mathbf v,S)$ is the fundamental object of RSVP field theory.

\begin{remark}
It is important not to confuse the entropy field $S$ with the classical thermodynamic entropy, although there are interpretative analogies.  In RSVP, $S$ is a geometric quantity arising from derived symplectic structures.
\end{remark}

\section{Derived Geometry}

The symbols $\mathbb{L}*\mathcal X$ and $\mathbb{T}*\mathcal X$ denote the cotangent and tangent complexes of a derived stack $\mathcal X$.  The shifted symplectic structure on $\Plenum$ is denoted by $\omega_{[-1]}$.

\section{Variational Notation}

The classical BV action is denoted $\mathcal{S}$, and the classical master equation is ${\mathcal S,\mathcal S}=0$.  The BV bracket is written ${-,-}_{BV}$ or simply ${-,-}$ when unambiguous.

\section{Constants and Parameters}

Physical constants such as $c$ and $G$ may be set to unity except in comparative formulas with observational cosmology.  Dimensional analysis is included in Volume~II as needed.

\section{End of Appendix D}

These conventions are used uniformly throughout all volumes of the monograph.  Any deviation in later chapters will be indicated explicitly.

